\documentclass[10.5pt,a5paper]{article}

\usepackage{fontspec}
\setmainfont{DejaVu Serif}
\setsansfont{DejaVu Sans}
\setmonofont{DejaVu Sans Mono}
\usepackage[english]{babel}
\usepackage{geometry}
\geometry{left=13mm,right=13mm,top=14mm,bottom=16mm}
\usepackage{microtype}
\usepackage{array,tabularx,booktabs}
\usepackage{enumitem}
\usepackage{xcolor}
\usepackage{hyperref}
\usepackage{fancyhdr}
\usepackage{titlesec}
\usepackage{tcolorbox}
\usepackage{qrcode}
\usepackage{graphicx}

% Build identity — same file as the passport and label. Fallbacks let this
% .tex compile standalone in a TeX editor without `make`.
\IfFileExists{build_info.tex}{\input{build_info}}{%
  \newcommand{\hwversion}{v2.1}%
  \newcommand{\fwversion}{dev}%
  \newcommand{\builddate}{unknown}%
  \newcommand{\buildcommit}{unknown}%
  \newcommand{\docrev}{unknown}%
}

% Public documentation site (GitHub Pages). Also the URL encoded in the
% QR code on the enclosure label — change in one place and everything follows.
\newcommand{\docsurl}{https://mrzlohex.github.io/dc-ups-2ch/}

\definecolor{accent}{HTML}{1E5B4F}
\definecolor{lightaccent}{HTML}{EAF3F1}
\definecolor{warn}{HTML}{8B1E1E}
\definecolor{lightwarn}{HTML}{F9ECEC}

\hypersetup{
  colorlinks=true,
  linkcolor=black,
  urlcolor=accent,
  pdfauthor={DC-UPS-2CH manufacturer},
  pdftitle={DC-UPS-2CH quick user manual \fwversion}
}

\setlength{\parindent}{0pt}
\setlength{\parskip}{0.4em}
\renewcommand{\arraystretch}{1.15}
\setlist[itemize]{leftmargin=5.5mm,itemsep=1.5pt,topsep=2pt}
\setlist[enumerate]{leftmargin=6.5mm,itemsep=2pt,topsep=2pt}

\titleformat{\section}{\large\bfseries\color{accent}}{\thesection.}{0.45em}{}
\titleformat{\subsection}{\normalsize\bfseries}{\thesubsection.}{0.45em}{}

\pagestyle{fancy}
\fancyhf{}
\lhead{\sffamily DC-UPS-2CH}
\rhead{\sffamily User manual}
\cfoot{\sffamily\small \thepage}
\renewcommand{\headrulewidth}{0.35pt}
\setlength{\headheight}{13pt}

\newtcolorbox{infobox}[1][]{colback=lightaccent,colframe=accent,boxrule=0.7pt,arc=1.5mm,left=2mm,right=2mm,top=1.5mm,bottom=1.5mm,#1}
\newtcolorbox{warnbox}[1][]{colback=lightwarn,colframe=warn,boxrule=0.7pt,arc=1.5mm,left=2mm,right=2mm,top=1.5mm,bottom=1.5mm,#1}

\begin{document}

\begin{titlepage}
\centering
\vspace*{10mm}
{\sffamily\bfseries\fontsize{24}{28}\selectfont DC-UPS-2CH\\[2mm]}
{\sffamily\Large Quick user manual\\[5mm]}
{\sffamily\small for firmware \fwversion\ (build \buildcommit, \builddate)}

\vspace{9mm}
\begin{infobox}
\centering
Dual-channel backup power supply for router, ONT and other compatible low-current network equipment.
\end{infobox}

\vspace{7mm}
\begin{tabularx}{\textwidth}{@{}lX@{}}
\toprule
Model & DC-UPS-2CH \hwversion \\
Serial number & \rule{45mm}{0.4pt} \\
Date of transfer & \rule{45mm}{0.4pt} \\
Output ROUTER & \rule{18mm}{0.4pt} V / \rule{18mm}{0.4pt} A \\
Output ONT & \rule{18mm}{0.4pt} V / \rule{18mm}{0.4pt} A \\
\bottomrule
\end{tabularx}

\vfill
\begin{warnbox}
\textbf{Warning.} Dangerous 230~V mains voltage is present inside the enclosure. The user must not open the enclosure while the mains or the battery is connected.
\end{warnbox}
\vfill
{\small Keep this manual next to the device.}
\end{titlepage}

\section{What the device does}

DC-UPS-2CH automatically maintains power on two independent outputs. Under normal conditions the device is powered from the mains and keeps the battery in a charged state. When 230~V mains is lost the switch-over to the battery happens automatically, with no user action required.

\begin{tabularx}{\textwidth}{@{}p{44mm}X@{}}
\toprule
\textbf{Parameter} & \textbf{Value} \\
\midrule
Mains input & 230~V AC, 50~Hz \\
Battery & 12~V, 7~Ah, AGM/SLA \\
Float-charge voltage & 13.8~V \\
Charge current limit & 1.5~A \\
Battery early warning & 11.5~V \\
Load disconnect (LVD) & 11.0~V \\
Load re-enable threshold & 12.5~V \\
\bottomrule
\end{tabularx}

\begin{infobox}[title=Normal operation]
Nothing needs to be switched over during a mains outage. The outputs keep running on the battery until its voltage reaches the protection threshold.
\end{infobox}

\section{Daily use}

\subsection{Turning on}
\begin{enumerate}
  \item Confirm the required outputs are connected.
  \item If the enclosure has separate mechanical output switches, set the required channels to the ON position.
  \item Apply 230~V mains.
  \item Wait for the connected equipment and Wi-Fi to come up.
\end{enumerate}

Under normal operation no further manual control is required.

\subsection{During a mains outage}
The device automatically switches to battery. If a network path is available, an ntfy notification is sent. At battery voltage 11.5~V an early-warning notification is sent. At 11.0~V both managed outputs are shut off to protect the battery.

After protective shut-off the ESP32 enters emergency deep sleep. It wakes periodically and checks whether the mains has returned. When it does, the device automatically resumes normal operation.

\begin{warnbox}[title=Do not confuse with Shelf Sleep]
After emergency deep sleep the device wakes automatically once the mains returns. In \textbf{Shelf Sleep} the return of 230~V does \textbf{not} wake the device.
\end{warnbox}

\section{Button and indication}

\subsection{Service button}
\begin{tabularx}{\textwidth}{@{}p{40mm}X@{}}
\toprule
\textbf{Gesture} & \textbf{Result} \\
\midrule
1 short press & Send full status message to ntfy \\
2 short presses & Open \texttt{DC-UPS-Setup} access point for about 15 minutes \\
Hold for 10~s & Enter Shelf Sleep \\
Any press while in Shelf Sleep & Wake the device and perform a normal start-up \\
\bottomrule
\end{tabularx}

\subsection{Button acknowledgement}
\begin{itemize}
  \item 2 green blinks --- status successfully sent to ntfy;
  \item 3 red blinks --- status could not be sent;
  \item short simultaneous flash of both LEDs --- setup AP has been opened;
  \item 3 red blinks before shutdown --- Shelf Sleep has been accepted.
\end{itemize}

\subsection{Normal indication}
\begin{tabularx}{\textwidth}{@{}p{35mm}X@{}}
\toprule
Green LED & external power / mains present \\
Red LED & battery state, warning or emergency mode \\
\bottomrule
\end{tabularx}

While the device is booting and joining Wi-Fi, the green LED blinks at about 2~Hz. When the setup access point is active, both LEDs alternate slowly (green/red, 500~ms each side).

\section{Shelf Sleep}

Shelf Sleep is intended for transport and long-term storage, when both outputs and Wi-Fi must be completely off.

\textbf{To power the device down for storage:}
\begin{enumerate}
  \item Hold the service button for about 10 seconds.
  \item Wait for three red blinks.
  \item Release the button.
  \item Both outputs will switch off, Wi-Fi will go down, the ESP32 will enter deep sleep with no periodic wake-up.
\end{enumerate}

\textbf{To power the device on again:} press the service button once.

\begin{warnbox}[title=Important]
In Shelf Sleep the device wakes \textbf{only from a physical button press}. Connecting or restoring 230~V mains does not wake it.
\end{warnbox}

\section{Web panel and Wi-Fi}

Once connected to the home network, the panel is available at:
\begin{center}
\large\texttt{http://dc-ups.local/}
\end{center}

The IP address delivered through ntfy on every successful Wi-Fi connect can also be used directly.

The panel shows voltages, mains / output / Wi-Fi / internet status and the event log; lets the user manage ROUTER/ONT modes; perform channel restarts; and change Wi-Fi, ntfy and protection settings.

\subsection{If the home Wi-Fi is not available}
Double-press the service button quickly. An access point will appear:

\begin{center}
\texttt{DC-UPS-Setup}
\end{center}

Default password: \texttt{dc-ups-setup}, unless it was changed during commissioning.

\section{ntfy}

The device sends notifications about important events on its own. In addition, one of the following commands may be posted to the same topic:

\begin{center}
\texttt{!ups}\qquad \texttt{!ups status}\qquad \texttt{!ups ping}
\end{center}

The commands only \textbf{query} state and \textbf{do not control the outputs}. The reply includes battery voltage, 24~V rail, mains status, ROUTER/ONT, internet, uptime, RSSI and IP.

\begin{infobox}
ntfy commands only work while the ESP32 is awake, connected to Wi-Fi and has internet access. In emergency deep sleep and in Shelf Sleep the Wi-Fi radio is off.
\end{infobox}

\section{Connecting equipment}

Before connecting a new router, ONT or other device, always verify:
\begin{itemize}
  \item required output voltage;
  \item barrel plug polarity;
  \item permissible current of the corresponding output;
  \item the ROUTER / ONT labelling on the enclosure.
\end{itemize}

For the standard cylindrical DC connector of this unit, polarity must match the marking on the enclosure. Do not connect equipment when the required voltage or polarity is unknown.

\section{Troubleshooting}

\begin{tabularx}{\textwidth}{@{}p{43mm}X@{}}
\toprule
\textbf{Symptom} & \textbf{What to do} \\
\midrule
No power on the outputs & Check mains, mechanical switches, battery voltage and the per-channel mode in the panel \\
Web panel not reachable & Try \texttt{dc-ups.local}; if it fails, double-press the button and connect to the setup AP \\
Outputs shut off after a mains outage & Normal LVD on low battery; wait for the mains to return \\
Device did not start after mains was restored & Shelf Sleep may be active --- press the service button \\
3 red blinks after a short press & Status was not sent to ntfy; check Wi-Fi / internet / topic \\
\bottomrule
\end{tabularx}

If the fault repeats, do not open the device while it is energised and contact the manufacturer or system integrator.

\section{Safety and maintenance}

\begin{itemize}
  \item do not operate the device with the enclosure open;
  \item do not allow water, condensation or airflow obstruction;
  \item do not short-circuit the battery terminals;
  \item do not exceed the output ratings printed on the label;
  \item the battery is a consumable part and loses capacity over time;
  \item if the autonomous runtime shortens noticeably, the battery should be checked and, if necessary, replaced according to the full service documentation.
\end{itemize}

\begin{warnbox}
Before opening the enclosure disconnect 230~V mains and the battery. Internal servicing is intended for trained personnel.
\end{warnbox}

\newpage
\section*{Quick reference}

\begin{infobox}
\textbf{Normal operation:} nothing to switch.\\
\textbf{Mains lost:} device works from the battery automatically.\\
\textbf{Mains returned:} device recovers from emergency deep sleep automatically.\\
\textbf{1x button:} status to ntfy.\\
\textbf{2x button:} setup AP.\\
\textbf{Hold 10~s:} Shelf Sleep.\\
\textbf{Wake from storage:} press the button.\\
\textbf{Panel:} \texttt{http://dc-ups.local/}.\\
\textbf{ntfy:} \texttt{!ups status} or \texttt{!ups ping}.
\end{infobox}

\vspace{6mm}
\begin{center}
\begin{tabular}{@{}c@{\hspace{14mm}}c@{}}
\qrcode[height=28mm]{http://dc-ups.local/} & \qrcode[height=28mm]{\docsurl} \\[1.5mm]
\\
\\
{\small\parbox{40mm}{\centering Local web panel\\(same LAN, mDNS)}} &
{\small\parbox{40mm}{\centering Online documentation\\(full passport, PDF)}} \\
\end{tabular}
\end{center}

\vfill
\begin{tabularx}{\textwidth}{@{}lX@{}}
Manufacturer contact: & \rule{50mm}{0.4pt} \\
Document date: & \builddate \\
Revision: & User manual \fwversion\ / \docrev \\
\end{tabularx}

\end{document}
